\section{Rational polytopes and normal lists}\label{sec:encoding}

The strict volume margin in the proof gives a polyhedral consequence of
Theorem~\ref{thm:main} without changing its constant.

\begin{corollary}\label{cor:rational-polytope}
With the same $c$ and $n_0$ as in Theorem~\ref{thm:main}, for every
$n\ge n_0$ there is a centrally symmetric polytope
$P_n\subset\mathbb R^n$ with rational vertices such that
\[
 \theta_T(P_n)>c n\log n.
\]
\end{corollary}
\begin{proof}
Fix the successful body $K$ in the proof of Theorem~\ref{thm:main},
and put $R=c n\log n$. It satisfies
$|K|>2v\omega_na^n$, and no periodic covering by $K$ has center
intensity at most $\rho_0=R/(2v\omega_na^n)$.

Enumerate the rational points of $\operatorname{int}K$ as
$q_1,q_2,\ldots$ and set
\[
 P_j=\operatorname{conv}\{\pm q_1,\ldots,\pm q_j\}.
\]
These are centrally symmetric rational polytopes contained in
$\operatorname{int}K$, and they increase with $j$. Their union is
$\operatorname{int}K$: every interior point lies inside a sufficiently
small simplex with rational vertices in $\operatorname{int}K$, and
all its vertices eventually occur in the enumeration. Continuity of
volume from below, together with the fact that the boundary of a convex
body has measure zero, gives $|P_j|\to|K|$. Choose $j$ so that
$|P_j|>2v\omega_na^n$; in particular, this $P_j$ has nonempty interior.

Every periodic covering by $P_j$ is a covering by $K$, because
$P_j\subset K$. Its center intensity therefore exceeds $\rho_0$.
Taking the infimum of the covering densities and using
Proposition~\ref{prop:periodization} yields
\[
 \theta_T(P_j)=\theta_{\mathrm{per}}(P_j)
       \ge\rho_0|P_j|>R.
\]
This argument uses both inclusion and the retained volume margin.
\end{proof}

A different finite description retains the spherical truncation $bB$
and makes all slab normals rational unit vectors. These bodies need not
be polytopes.

\begin{proposition}\label{prop:rational}
The bodies in Theorem~\ref{thm:main} can be chosen in the form
\[
 K=bB\cap\bigcap_{j=1}^m\{x:|u_j\cdot x|\le1\},
 \qquad u_j\in\mathbb Q^n\cap S^{n-1},
\]
with $m<\infty$.
\end{proposition}
\begin{proof}
Fix a successful finite normal list $U=(u_1,\ldots,u_m)$ from the proof of
Theorem~\ref{thm:main}. The class $\mathcal F_n$ is finite and independent
of $U$. For each $P\in\mathcal F_n$, select a point $y_P\in B_D$ outside
all of its translates of $K^+(U)$. For every center $p$ in that list,
either $|y_P-p|>b+3\eta$, or one of the normals satisfies
\[
 |u_j\cdot(y_P-p)|>1+3\eta.
\]
Each required inequality is strict. There are finitely many patterns,
centers, and chosen certificates, so all these inequalities persist under
a sufficiently small simultaneous perturbation of the normals. Empty
patterns need no certificates.

The strict volume test persists as well. If $U'$ is a list of unit normals
with $|u_j-u'_j|\le\delta$ for every $j$, then
\begin{equation}\label{eq:volume-continuity}
 (1+b\delta)^{-1}K(U)\subset K(U')\subset(1+b\delta)K(U).
\end{equation}
Indeed, if $x\in K(U)$, then $|x|\le b$ and
$|u'_j\cdot x|\le1+b\delta$; scaling by $(1+b\delta)^{-1}$ proves the
first inclusion. Interchanging $U$ and $U'$ proves the second. Volumes
therefore converge as $\delta\to0$.

Rational unit vectors are dense in $S^{n-1}$ for $n\ge2$, by the
stereographic parametrization
\[
 t\longmapsto\frac{(2t,1-|t|^2)}{1+|t|^2},
 \qquad t\in\mathbb Q^{n-1},
\]
and approximation of its omitted pole. Approximate all normals closely
enough to preserve the holes and the strict volume test. The localization
argument applies to every finite normal list, so the resulting rational
list yields the same lower bound.

\end{proof}

One can specify a choice exactly by ordering the finite rational normal
lists. Write each coordinate in reduced form with positive denominator,
and let its height be the maximum of the absolute numerator and denominator.
Order lists by their length plus maximum coordinate height (height $1$
for the empty list), breaking each finite tie lexicographically.
The lists satisfying the strict volume test and every finite-pattern
noncoverage test form a nonempty set, so they have a first member.
This specifies a finite mathematical object; it does not assert an
algorithm for deciding the volume comparisons or an effective bound for
the first admissible dimension.
